\documentclass[10pt]{article}

\usepackage[utf8]{inputenc}

\usepackage[T1]{fontenc}

\usepackage{amsmath}

\usepackage{amsfonts}

\usepackage{amssymb}

\usepackage[version=4]{mhchem}

\usepackage{stmaryrd}

\usepackage{bbold}

\usepackage{graphicx}

\usepackage[export]{adjustbox}

\graphicspath{ {./images/} }

\usepackage{mathrsfs}



\begin{document}

щим изолированных точек). Другими словами, $F$ совІадает со своим производны.и множеством. Примеры С. м.: $\mathbb{R}^{n}, \mathbb{C}^{n}$, канторово множество. М. И. Войцеховский. СОВЕРІЕННОЕ НЕПРИВОДИМОЕ ОТОБРАЖЕНИЕ - совершенное отображение $f$ пространства $X$ на пространство $Y$, являющееся неприводимым (т. е. $Y$ не является образом никакого замкнутого в $X$ множества, отличного от $X$ ). $\quad M$. И. Войцеховский.



COBEРIIЕННОЕ ОТОБРАЖЕНИЕ - непрерывное замкнутое отображение топологич. пространств, при к-ром прообразы всех точек бикомпактны. С. о. во многом аналогичны непрерывным отображениям бикомпактов в хаусдорфовы пространства (каждое такое отображение совершенно), но сферой действия имеют класс всех топологич. пространств. В классе вполне регулярных пространств С. о. характеризуются существованием у них непрерывного продолжения на нек-рые бикомпактные распирения, при к-ром наросты распирений отображаются в наросты. С. о. сохраняет метризуемость, паракомпактность, вес, полноту по Чеху в сторону образа; другие инварианты (напр., характер пространства) оно преобразует правильным образом. Класс С. о. замкнут относительно операций произведения и композиции. Сужение С. о. на замкнутое подпространство является С. о. (не так обстоит дело для факторных и открытых отображений).



Названные свойства С. о. привели к тому, что класс этих отображений стал играть стержневую роль в классификации топологич. Іоространств. Прообразы метрич. пространств при С. о. охарактеризованы как паракомпактные перистые $(p)$-гространства. Класс паракомпактных $p$-пространств замкнут уже в обе стороны относительно С. о. Важным свойством С. о. является возможность сузить каждое из них на нек-рое замкнутое подпространство, не уменьшая образа, так, чтобы получившееся отображение было неприводимым - не допускало дальнейшего сужения на замкнутое подпространство без уменьшения образа. Неприводимые С. о. являются отправной точкой построения теории абсолютов топологич. пространств. При неприводимом С. $о$.



\includegraphics[max width=\textwidth, center]{2023_07_09_a1c99833d778c0f4e201g-1(2)}

равства и тисло Суслина образа равно числу Суслина отображаемого пространства. Если вполне регулярное $T_{1}$-пространство $X$ отображается на вшолне регулярное $T_{1}$-пространство $Y$ посредством С. о., то $X$ гомеоморфно замкнутому подпространству топологич. произведения пространства $Y$ на нек-рый бикомпакт. Диагональное произведение С. о. и непрерывного отображения всегда является С. о., в частности диагональное произведение С. о. и уплотнения является гомеоморфизмом, и если топологическое пространств-о совершенно отображается и уілотняется на нек-рое (вообще говоря, другое) метрическое пространство, то оно само метризуемо.



Лum.: [1] А р х ангельс ікй̈ А. В., Пономаp е в В. И. О Оновы общей топологии в задачах и упражнениях,

M., $1974 ;[2]$ Б у р б а к и Н., Общая топология. Основные структуры, пер. с франц., м., 1968.



СОВЕРІІЕННЕ ПОЛЕ - поле $k$, Любой многочлен над к-рым сепарабелен. Иначе говоря, любое алгебраич. расширение поля $k$ - сепарабельное расширение. Все остальные поля наз. н е с о в е р ш е н н ы м. Все поля характеристики 0 совершенны. Поле $k$ конечной характеристики $p$ совершенно тогда и только тогда, когда $k=k p$, т. е. возведение в степень $p$ является автоморфизмом поля $k$. Конечные поля и алгебраическі замкнутые поля совершенны. Пример несовершенного поля - поле $F_{q}(X)$ рациональных функций над полем $F_{q}$, где $F_{q}$ - поле из $q=p^{n}$ элементов. С. П. $k$ совпадает с полем инвариантов группы всех $k$-автоморфизмов алгебраич. замыкания $\bar{k}$ поля $k$. Любое алгебраич. расширение С. 1. снова совершенно. Для произвольного поля $k$ характеристики $p>0$ с алгебраич. замыканием $\bar{k}$ поле



\[

k^{p-\infty}=\bigcup_{n} k^{p-n} \subset \bar{k}

\]



является наименышим С. п., содержащим $k$. Оно наз. с ов е р ше н ным замыка ни е м поля $k$ в $\bar{k}$. Лит.: [1] Б у р б а к и Н., Алгебра. Многочлены и поля. с к и й О., С а м ю э л ь П., Коммутативная алгебра, пер. с



\begin{center}

\includegraphics[max width=\textwidth]{2023_07_09_a1c99833d778c0f4e201g-1}

\end{center}



СО'ВЕРIIIEННОЕ ЧИСЛО - целое положительное число, обладающее свойством, что оно совпадает с суммой всех своих положитөльных делителей, отличных от самого этого числа.



Таким образом, целое число $n \geqslant 1$ является С. ч., если



\[

n=\sum_{0<d<n, d / n} 1

\]



С. ч. являются, напр., числа $6,28,496,8128,33550336, \ldots$



С. ч. тесно связаны с простыми Мерсенна числами, т. е. с простыми числами вида $2^{m}-1$. Еще Евклид установил, что число



\[

n=2^{m-1}\left(2^{m}-1\right)

\]



является совершснным, ссли $2^{m}-1$ - простое число. J. Эйлер (L. Euler) показал, что этими числами исчерпываются все четные С. ч.



До сих пор (1983) неизвестно, будет ли конечным или бесконечным множество четных С. ч., т. е. неизвестно, будет ли конечным или бесконечным множество простых чисел Мерсенна $2^{m}-1$. Неизвестно также, существуют ли нечетные С. ч.



До 1983 найдено 27 четных С. ч. Первые 23 из них соответствуют следующим звачениям $m: 2,3,5,7,13$, $17,19,31,61,89,107,127,521,607,1279,2203,2281$, $3217,4219,4423,9689,9941,11213$. Список С. ч. с 12-го по 24-е указан в [2]. Четные С. ч. с 25-го по 27-е соответствуют следующим значениям $m: 21701,23209,44497$ (см. [3]). Показано (см. [4]), тто нечетных С. ч. нет в интервале от 1 до $10^{50}$.



Лum.: [1] D i c k s o n L. E., History of the theory of numbers, v. 1, Wash., 1919, repr., N.- Y., 1952; [2] N a n k a r M. L., "Ganita Bhāratī", 1979, v. 1, N 1-2, p. 7-8; [3] Sl ow i n s k i D., "M. Recreational Math.", $1979, \mathrm{v} .11, \mathrm{p} .258-61 ;$ [. A. Cтепанов.



СОВМЕСТНОЕ РАСПРЕДЕЛЕНИЕ - общиЙ термин, относящийся к распределению нескольких случайных величин, заданных на одном и том же вероятностном пространстве. Пусть случайные величины $X_{1}, \ldots, X_{n}$ определены на вероятностном пространстве $\{\Omega, \mathcal{A}, \mathrm{P}\}$ и принимают значения в измеримых простравствах $\left(\mathfrak{X}_{k}, \mathscr{B}_{k}\right)$. С о в м е с т н ы м распределением этих величин наз. функция $P_{X_{1} \ldots X_{n}}\left(B_{1}, \ldots, B_{n}\right)$, определенная на множествах $B_{1} \in \mathscr{B}_{1}, \ldots, B_{n} \in \mathscr{B}_{n}$ как



\[

P_{X_{1} \ldots X_{n}}\left(B_{1}, \ldots, B_{n}\right)=\mathrm{P}\left\{X_{1} \in B_{1}, \ldots, X_{n} \in B_{n}\right\} .

\]



В связи с C. р. говорят о с о в м е с т н о й ф ун кции распределения и ос ов ме с тной пл лотности вероятности.



Если $X_{1}, \ldots, X_{n}$ - обычные действительные случайные величины, то C. p. есть распределение случайного вектора $\left(X_{1}, \ldots, X_{n}\right)$ в пространстве $\mathbb{R} n$ (cм. $M_{\text {ного- }}$ мерное распределение). Если $X(t), t \in T,-$ случайный процесс, то С. р. значений $X\left(t_{1}\right), \ldots, \bar{X}\left(t_{n}\right)$ при $t_{1}, \ldots$, $t_{n} \in T$ наз. кон е чном е рным и р а сп р е д е-



\includegraphics[max width=\textwidth, center]{2023_07_09_a1c99833d778c0f4e201g-1(1)}

Лum.: [1] П $\mathrm{p}$ о х о р о в Ю. В., Р о а н о в Ю.. А., Тео-

рия вероятностей, 2 изд., М., 1973. СОВМЕСТНОСТЬ МЕТОДОВ СУММИРОВАНИЯ свойство методов суммирования, состоящее в непротиворечивости результатов применения этих методов. Методы $A$ и $B$ совместны, если они не могут суммировать одну и ту же последовательность или ряд к различным пределам, в противном случае они наз. н е с о в-





\end{document}